\documentclass[addpoints, 12pt]{exam}
\usepackage[T1]{fontenc}
\usepackage[utf8]{inputenc}
\usepackage[brazil]{babel}
\usepackage{graphicx}
\usepackage{geometry}
\usepackage[export]{adjustbox} 
\usepackage{tikz}
\usepackage{tabularx} 
\usepackage{multicol} 
\usepackage{amsmath, amssymb}
\usepackage[shortlabels]{enumitem}        

% --- CONFIGURAÇÕES DAS COLUNAS (O FIM DO EFEITO CHICLETE) ---
\setlength{\columnsep}{1cm}       
\setlength{\columnseprule}{0.4pt} 
\raggedcolumns % 👈 MÁGICA 1: Proíbe o LaTeX de esticar as colunas no meio da prova!

% 👈 MÁGICA 2: Corta a gordura (espaços enormes) entre as alternativas (A, B, C, D)
\renewcommand{\choiceshook}{
  \setlength{\itemsep}{0pt}
  \setlength{\topsep}{0.1cm}
}
% 👈 MÁGICA 3: Mantém as questões levemente separadas, mas sem exageros
\renewcommand{\questionshook}{
  \setlength{\itemsep}{0.3cm}
}

% MARGENS
\geometry{left=1cm, right=1cm, top=1cm, bottom=2cm}

% 👇 COLE ESTE BLOCO AQUI 👇
% --- CONFIGURAÇÕES DE RODAPÉ (CÓDIGO SECRETO) ---
\pagestyle{foot}
\firstpagefooter{}{\small \texttt{Código de Verificação: << codigo_secreto >>}}{Pág. \thepage}
\runningfooter{}{\small \texttt{Código de Verificação: << codigo_secreto >>}}{Pág. \thepage}

\begin{document}

<% for _ in range(num_copias) %>
% --- 1. CABEÇALHO ESTILO ENGENHARIA ---
\noindent
% --- BLOCO 1: LOGO (18% da largura) ---
\begin{minipage}[t]{0.18\textwidth}
    \vspace{0pt}
    \fbox{\begin{minipage}[c][3.2cm][c]{\dimexpr\linewidth-2\fboxsep-2\fboxrule}
    \centering
        \includegraphics[width=0.95\linewidth, height=3cm, keepaspectratio]{<< logo_path >>}
    \end{minipage}}
\end{minipage}%
\hfill
% --- BLOCO 2: TEXTOS (58% da largura) ---
\begin{minipage}[t]{0.58\textwidth}
    \vspace{0pt} 
    {\large \textbf{<< instituicao >>}} \\[0.15cm]
    Professor(a): << professor_nome >> \\[0.1cm]
    Disciplina: << disciplina_nome >> \hrulefill \\[0.1cm]
    Curso: << curso >> \hrulefill \\[0.1cm]
    <% if aluno_nome %>
        \textbf{Aluno(a):} << aluno_nome >> \hrulefill \\[0.1cm]
        \textbf{RA:} << aluno_ra >> \hfill Turma: << turma >> \rule{4cm}{0.1pt}
    <% else %>
        Aluno(a): \hrulefill \\[0.1cm]
        RA: \rule{4cm}{0.1pt} \hfill Turma: << turma >> \rule{4cm}{0.1pt}
    <% endif %>
\end{minipage}%
\hfill
% --- BLOCO 3: NOTA e DATA (20% da largura) ---
\begin{minipage}[t]{0.20\textwidth}
    \vspace{0pt}
    % Caixa da nota limpa para preenchimento manual
    \fbox{\begin{minipage}[t][2.5cm][t]{\dimexpr\linewidth-2\fboxsep-2\fboxrule}
        \centering \vspace{0.2cm}
        \textbf{Nota}
    \end{minipage}}
    \vspace{0.2cm} \\
    Data: << data >>\hrulefill \\[0.1cm]
    
\end{minipage}

\vspace{0.4cm}
% --- TÍTULO EM DESTAQUE ---
\vspace{0.4cm}
\begin{center}
    {\Huge \textbf{<< titulo_documento >>}} % Título bem grande e negrito
    \\[0.2cm] \hrule % 👈 Adiciona uma linha horizontal elegante
\end{center}
\vspace{0.2cm}
% --------------------------

% --- INSTRUÇÕES ---
\noindent\fbox{\begin{minipage}{\dimexpr\linewidth-2\fboxsep-2\fboxrule}
    \textbf{Instruções:}\\[0.2em]
    << instrucoes_texto >>
\end{minipage}}

\vspace{0.5cm}

% --- 2. QUESTÕES ---
<% if colunas == 2 %>
\begin{multicols*}{2}
<% endif %>
\begin{questions}
<% for q in questoes %>
    \question \textbf{(<< q.pontos >> pontos)} << q.enunciado >>

    <% if q.imagem %>
    \par\vspace{0.3cm}
    {\centering \includegraphics[max width=0.95\linewidth, max height=8cm, keepaspectratio]{<< q.imagem >>} \par}
    \vspace{0.1cm}
    <% endif %>

    <% if q.alternativas %>
    \begin{choices}
        <% for alt in q.alternativas %>
        \choice << alt.texto >>
        <% if alt.imagem %>
        \par \includegraphics[max width=0.85\linewidth, max height=4cm, keepaspectratio]{<< alt.imagem >>}
        <% endif %>
        <% endfor %>
    \end{choices}
    <% endif %>

    <% if q.espaco != "Nenhum" %>
        \vspace{0.2cm}
        <% if q.espaco == "Linhas" and q.num_linhas > 0 %>
            % 👇 A VERDADEIRA MÁGICA: Linhas independentes que o multicol consegue fatiar 👇
            <% for i in range(q.num_linhas) %>
                \noindent\rule{\linewidth}{0.4pt}\par\vspace{0.55cm}
            <% endfor %>
        <% elif q.espaco == "Quadriculado" %>
            \fillwithgrid{<< q.espaco_linhas >>cm}
        <% elif q.espaco == "Caixa Vazia" %>
            \makeemptybox{<< q.espaco_linhas >>cm}
        <% endif %>
    <% endif %>
    \vspace{0.4cm}
<% endfor %>
\end{questions}
<% if colunas == 2 %>
\end{multicols*}
<% endif %>

% --- 3. CARTÃO RESPOSTA (Otimizado para OpenCV) ---
\newpage
\begin{center} \textbf{\huge CARTÃO RESPOSTA OFICIAL} \end{center}

% 4 Quadrados pretos nos cantos
\begin{tikzpicture}[remember picture, overlay]
    \node at ([shift={(1.5cm,-1.5cm)}]current page.north west) {\rule{0.8cm}{0.8cm}};
    \node at ([shift={(-1.5cm,-1.5cm)}]current page.north east) {\rule{0.8cm}{0.8cm}};
    \node at ([shift={(1.5cm,1.5cm)}]current page.south west) {\rule{0.8cm}{0.8cm}};
    \node at ([shift={(-1.5cm,1.5cm)}]current page.south east) {\rule{0.8cm}{0.8cm}};
\end{tikzpicture}

\vspace{0.5cm}

% CAIXA DE IDENTIFICAÇÃO (LIMPA E SEM VERSÃO DA PROVA)
\begin{center}
\begin{tabular}{|p{0.45\textwidth}|p{0.5\textwidth}|}
    \hline
    \begin{center}
        \vspace{0.2cm}
        \includegraphics[width=3.5cm]{<< qr_path >>}
        \vspace{0.2cm}
    \end{center} &
    
    \begin{center}
        \vspace{0.2cm}
        \textbf{\large IDENTIFICAÇÃO DO ALUNO}
        
        \vspace{0.8cm}
        \textbf{\footnotesize NOME:} \\
        \large << aluno_nome >>
        
        \vspace{0.8cm}
        \textbf{\footnotesize RA:} \\
        \large << aluno_ra >>
        \vspace{0.6cm}
    \end{center} \\
    \hline
\end{tabular}
\end{center}

\vspace{1cm}
\begin{center}
\fbox{
\begin{minipage}{0.8\textwidth}
    \centering\fontsize{8}{10}\sffamily
    \textbf{Instruções para Questões Numéricas:} \\
    O primeiro bloco de bolinhas representa a \textbf{Dezena} e o segundo a \textbf{Unidade}. \\
    Exemplo: Para a resposta \textbf{42}, marque o (4) no primeiro bloco e o (2) no segundo.
\end{minipage}
}
\end{center}

\begin{center}
    \textbf{\Large RESPOSTAS OBJETIVAS} \\ \vspace{0.4cm}
    \renewcommand{\arraystretch}{2} % Aumenta bem o espaço entre as linhas
    
   % TABELA DAS BOLINHAS (ALINHAMENTO MATEMÁTICO PERFEITO)
    \begin{tabular}{c c l}
        <% for q in questoes %>
        <% if q.tipo in ['Múltipla Escolha', 'Verdadeiro ou Falso', 'Numérica'] %>
            $\vcenter{\hbox{\textbf{\large << loop.index >>.}}}$ & 
            $\vcenter{\hbox{\rule{0.25cm}{0.25cm}}}$ & 
            $\vcenter{\hbox{
            <% if q.tipo == 'Numérica' %>
                \begin{tikzpicture}
                    % Bloco da DEZENA
                    \foreach \n in {0,...,9} { 
                        \draw[black!50, thin] (0.45*\n, 0) circle (0.16) node[text=black] {\sffamily\fontsize{5}{6}\selectfont \n}; 
                    }
                    % Ponto separador
                    \node at (4.4, 0) {\small $\cdot$};
                    % Bloco da UNIDADE
                    \foreach \n in {0,...,9} { 
                        \draw[black!50, thin] (4.7 + 0.45*\n, 0) circle (0.16) node[text=black] {\sffamily\fontsize{5}{6}\selectfont \n}; 
                    }
                \end{tikzpicture}
            <% elif q.tipo == 'Verdadeiro ou Falso' %>
                \begin{tikzpicture}
                    \draw[black!50, thin] (0, 0) circle (0.16) node[text=black] {\sffamily\fontsize{5}{6}\selectfont V};
                    \draw[black!50, thin] (0.45, 0) circle (0.16) node[text=black] {\sffamily\fontsize{5}{6}\selectfont F};
                \end{tikzpicture}
            <% else %> % Múltipla Escolha
                \begin{tikzpicture}
                    \foreach \l [count=\i from 0] in {A,B,C,D,E} {
                        \draw[black!50, thin] (0.45*\i, 0) circle (0.16) node[text=black] {\sffamily\fontsize{5}{6}\selectfont \l};
                    }
                \end{tikzpicture}
            <% endif %>
            }}$ \\
        <% endif %>
        <% endfor %>
    \end{tabular}
\end{center}

<% if not loop.last %>\newpage<% endif %>
<% endfor %>
\end{document}